\documentclass[12pt,a4paper]{article}

% -------------------------------------------------
% Sprache / Kodierung
% -------------------------------------------------
\usepackage[T1]{fontenc}
\usepackage[utf8]{inputenc}
\usepackage[ngerman]{babel}
\usepackage{lmodern}
\usepackage{microtype}

% -------------------------------------------------
% Mathematik
% -------------------------------------------------
\usepackage{amsmath,amssymb,amsthm,mathtools}

% -------------------------------------------------
% Layout / Grafik / Farben
% -------------------------------------------------
\usepackage[a4paper,margin=2.3cm]{geometry}
\usepackage{booktabs}
\usepackage{xcolor}
\usepackage[hidelinks,unicode]{hyperref}
\pdfstringdefDisableCommands{%
  \def\ü{ue}\def\ö{oe}\def\ä{ae}\def\Ü{Ue}\def\Ö{Oe}\def\Ä{Ae}\def\ss{ss}%
}
\usepackage[most]{tcolorbox}
\usepackage{enumitem}
\usepackage{array}

% -------------------------------------------------
% Farben (konsistent mit Vorgängerdokumenten)
% -------------------------------------------------
\definecolor{lückenrot}{RGB}{180,40,40}
\definecolor{bewiesengrün}{RGB}{30,100,50}
\definecolor{warnorange}{RGB}{200,100,10}
\definecolor{hintergrundblau}{RGB}{235,242,250}
\definecolor{hintergrundgelb}{RGB}{255,252,225}
\definecolor{adischviolett}{RGB}{90,20,140}
\definecolor{fehlerrot}{RGB}{200,0,0}
\definecolor{konjviolett}{RGB}{80,0,130}
\definecolor{e8grau}{RGB}{55,55,90}
\definecolor{kingmanblau}{RGB}{20,60,130}

% -------------------------------------------------
% tcolorbox-Stile
% -------------------------------------------------
\tcbset{
  lücke/.style={colback=red!8, colframe=lückenrot, fonttitle=\bfseries,
    title={Offene Lücke}},
  ergebnis/.style={colback=green!7, colframe=bewiesengrün,
    fonttitle=\bfseries, title={Bewiesenes Ergebnis}},
  warnung/.style={colback=orange!10, colframe=warnorange,
    fonttitle=\bfseries, title={Achtung / Heuristik}},
  adisch/.style={colback=violet!6, colframe=adischviolett,
    fonttitle=\bfseries, title={2-adische Perspektive}},
  kernidee/.style={colback=hintergrundblau, colframe=blue!60!black,
    fonttitle=\bfseries, title={Kernidee}},
  konj/.style={colback=violet!5, colframe=konjviolett, fonttitle=\bfseries,
    title={Konjektur (unbewiesen)}},
  lte/.style={colback=yellow!8, colframe=e8grau, fonttitle=\bfseries,
    title={LTE-Analyse}},
  hauptsatz/.style={colback=green!10, colframe=bewiesengrün!80!black,
    fonttitle=\bfseries\large, title={Satz (bewiesen)}},
  negativ/.style={colback=red!5, colframe=lückenrot!80!black,
    fonttitle=\bfseries, title={Negatives Ergebnis (bewiesen)}},
  kingman/.style={colback=blue!5, colframe=kingmanblau,
    fonttitle=\bfseries, title={Kingman-Perspektive}},
  tao/.style={colback=teal!5, colframe=teal!60!black,
    fonttitle=\bfseries, title={Verbindung zu Tao (2019)}},
}

% -------------------------------------------------
% Theorem-Umgebungen
% -------------------------------------------------
\newtheorem{definition}{Definition}[section]
\newtheorem{proposition}[definition]{Proposition}
\newtheorem{theorem}[definition]{Theorem}
\newtheorem{lemma}[definition]{Lemma}
\newtheorem{korollar}[definition]{Korollar}
\newtheorem{bemerkung}[definition]{Bemerkung}
\newtheorem{hypothese}[definition]{Hypothese}
\newtheorem{satz}[definition]{Satz}

\newenvironment{beweis}{\begin{proof}[Beweis]}{\end{proof}}

% -------------------------------------------------
% Makros
% -------------------------------------------------
\newcommand{\vii}{\nu_2}
\newcommand{\U}{\mathcal{U}}
\newcommand{\N}{\mathbb{N}}
\newcommand{\Z}{\mathbb{Z}}
\newcommand{\R}{\mathbb{R}}
\newcommand{\Ztwo}{\mathbb{Z}_2}
\newcommand{\abs}[1]{\lvert#1\rvert}
\newcommand{\atwo}[1]{\lvert#1\rvert_2}
\newcommand{\Epi}{\mathbb{E}_\pi}
\newcommand{\Lam}{\Lambda}

% -------------------------------------------------
% Dokument
% -------------------------------------------------
\title{%
  Kingmans Unteradditivitätssatz und die Collatz-Dynamik\\[0.3em]
  \large Birkhoff, Kingman und Ruelle-Perron-Frobenius im Vergleich ---
  welcher Satz passt?\\[0.5em]
  \normalsize Ergänzung zu \textit{collatz\_lte\_dichte.tex} und
  \textit{collatz\_schlussartikel.tex}%
}
\author{Thomas Hoffbauer}
\date{16.\ Juni 2026}

\begin{document}
\maketitle

\begin{abstract}
Wir analysieren, welcher Ergodensatz für den Collatz-Lyapunov-Exponenten
$\Lambda \approx -0{,}830$ am direktesten anwendbar ist.

\textbf{Hauptergebnis:} Die Block-Faktor-Folge ist \emph{additiv} (nicht
nur subadditiv), weshalb der \textbf{Birkhoff-Ergodensatz} und der
\textbf{Ruelle-Perron-Frobenius-Satz} (für die mod-12-Markov-Kette)
direkter anwendbar sind als Kingmans Satz. Kingman ist dennoch nützlich
für die \emph{subadditive Hülle} der LTE-korrigierten Folge.

\textbf{Verbindung zu Tao (2019):} Unser $\Lambda \approx -0{,}830$
quantifiziert explizit, was Tao mit Log-Dichte-$1$-Methoden implizit
zeigt: Fast alle Collatz-Bahnen haben einen negativen Lyapunov-Exponenten.

\textbf{Ehrliche Einschätzung:} Die Ergodizität von $T$ bleibt das
zentrale offene Problem. Birkhoff und Kingman liefern Konvergenz nur
unter dieser Voraussetzung --- und Ergodizität ist zur Collatz-Vermutung
äquivalent (über 2-adische Orbits).
\end{abstract}

\tableofcontents

\bigskip
\begin{tcolorbox}[warnung, title={Einordnung dieses Dokuments}]
Dieses Dokument enthält eine Kombination aus \textbf{bewiesenen Sätzen}
der Ergodentheorie (grün), \textbf{Anwendungsanalysen} auf Collatz
(orange, teils offen) und \textbf{Verbindungen zur Literatur}.
Die Collatz-Vermutung bleibt unbewiesen.
\end{tcolorbox}

% ====================================================
\section{Kingmans Unteradditivitätssatz}
\label{sec:kingman-satz}
% ====================================================

\subsection{Präzises Statement}

\begin{tcolorbox}[hauptsatz, title={Kingmans Unteradditivitätssatz (Kingman 1968)}]
\label{satz:kingman}
Sei $(X, \mathcal{F}, \mu, T)$ ein maßerhaltendes ergodisches System.
Sei $(f_n)_{n \geq 1}$ eine Folge von $L^1(\mu)$-Funktionen mit der
\emph{subadditiven} Eigenschaft:
\[
  f_{m+n}(x) \leq f_m(x) + f_n(T^m x)
  \quad \mu\text{-fast überall für alle } m,n \geq 1.
\]
Dann gilt:
\begin{enumerate}[label=(\roman*)]
  \item Die Folge $\frac{1}{n} f_n(x)$ konvergiert $\mu$-fast überall und
    in $L^1$ gegen eine $T$-invariante Funktion $f^* \in L^1(\mu)$.
  \item Für ergodisches $T$ ist $f^*$ $\mu$-fast überall konstant:
    \[
      f^*(x) = \inf_{n \geq 1} \frac{1}{n}\int f_n \,d\mu
              = \lim_{n \to \infty} \frac{1}{n}\int f_n \,d\mu
              \quad \mu\text{-f.ü.}
    \]
  \item Es gilt stets $\frac{1}{n}\int f_n \,d\mu \searrow f^*$ (monoton
    fallend, da $\frac{1}{n}\int f_n \leq \frac{1}{m}\int f_m$ für $m \mid n$
    aus Subadditivität folgt).
\end{enumerate}
\end{tcolorbox}

\begin{bemerkung}[Spezialisierung]
Für \emph{additive} Folgen, d.h.\ $f_n(x) = \sum_{k=0}^{n-1} \phi(T^k x)$,
ist die subadditive Ungleichung mit Gleichheit erfüllt. In diesem Fall
liefert Kingmans Satz nichts Neues gegenüber dem klassischen
Birkhoff-Ergodensatz. Kingman ist eine \emph{echte Verallgemeinerung}
nur für genuint subadditive Folgen, bei denen Ungleichheit herrscht.
\end{bemerkung}

\subsection{Superadditive Variante}

Für eine \emph{superadditive} Folge $(g_n)$, also
$g_{m+n}(x) \geq g_m(x) + g_n(T^m x)$, gilt das analoge Resultat
mit $\sup$ statt $\inf$. Das ist relevant für die subadditive Hülle
in §\,\ref{sec:kingman-lte}.

% ====================================================
\section{Das Collatz-System als ergodischer Rahmen}
\label{sec:system}
% ====================================================

\subsection{Die 2-adische Erweiterung}

Die klassische Collatz-Abbildung auf $\N_{\text{odd}} = 2\N+1$ ist keine
Selbstabbildung eines kompakten Raums. Die kanonische Einbettung nutzt den
Ring der 2-adischen ganzen Zahlen:

\begin{definition}[2-adischer Collatz-Raum]
Sei $\Ztwo$ der Ring der 2-adischen ganzen Zahlen mit ultrametrischer Topologie,
$\Ztwo^* = \{x \in \Ztwo : \atwo{x} = 1\}$ die Einheitengruppe
(``ungerade 2-adische Zahlen''). Die Collatz-Abbildung
\[
  T\colon \Ztwo^* \to \Ztwo^*,
  \qquad
  T(x) = \frac{3x+1}{2^{\vii(3x+1)}}
\]
ist auf $\Ztwo^*$ wohldefiniert und \emph{stetig} bezüglich der 2-adischen Topologie.
\end{definition}

Das Haar-Maß $\mu_H$ auf $\Ztwo^*$ ist das natürliche Wahrscheinlichkeitsmaß
(normiert: $\mu_H(\Ztwo^*) = 1$). Es ist unter $T$ \emph{nicht erhalten}, aber
quasi-invariant mit kontrollierter Radon-Nikodym-Ableitung.

\begin{tcolorbox}[adisch, title={Tao (2019): Quasi-Invarianz}]
Tao zeigt, dass für jedes $\varepsilon > 0$ die Menge
\[
  \bigl\{x \in \Ztwo^* : \sup_n T^n(x) > f(x)\bigr\}
\]
für geeignet wachsendes $f$ Haar-Maß null hat. Das impliziert, dass
$T$ kein $\mu_H$-erhaltendes System ist, aber ein \emph{dissipativer}
Endomorphismus mit messbarer Struktur.
\end{tcolorbox}

\subsection{Die Block-Zerlegung}

Auf der Ebene der \emph{Blöcke} ergibt sich ein besseres ergodisches Bild.
Jede ungerade Zahl $n$ wird nach ihrer EABC-Klasse (mod\,12) klassifiziert:

\begin{center}
\renewcommand{\arraystretch}{1.2}
\begin{tabular}{ccc}
\toprule
Klasse & $n \bmod 12$ & Block-Typ \\
\midrule
E & $1$ & reine E-Schritte, starke Kontraktion \\
A & $5$ & A-Schritt, moderate Kontraktion \\
B & $7$ & B-Schritt, schwache Kontraktion \\
C & $11$ & C-Kette gefolgt von A-Schritt \\
\bottomrule
\end{tabular}
\end{center}

Die Block-Iteration $T^{(1)}: \Ztwo^* \to \Ztwo^*$ wirft jeweils einen
Block ab und liefert den nächsten Startzustand. Die Zustandsraumreduktion
auf $\{E, A, B, C\} \cong \Z/12\Z$ ergibt eine \emph{endliche Markov-Kette}.

% ====================================================
\section{Birkhoff, Kingman, Ruelle-Perron-Frobenius — Vergleich}
\label{sec:vergleich}
% ====================================================

Für das Collatz-Block-System stehen drei Werkzeuge zur Verfügung.
Wir analysieren Stärken und Schwächen.

\subsection{Birkhoffs Ergodensatz}

\begin{theorem}[Birkhoff 1931]
Sei $(X, \mathcal{F}, \mu, T)$ ergodisch und $\phi \in L^1(\mu)$.
Dann gilt für $\mu$-fast alle $x$:
\[
  \frac{1}{n}\sum_{k=0}^{n-1}\phi(T^k x) \xrightarrow{n\to\infty} \int \phi \,d\mu.
\]
\end{theorem}

\textbf{Anwendung auf Collatz:} Definiere $\phi(x) = \log_2 F(B(x))$, wobei
$F(B)$ der Wachstumsfaktor des ersten Blocks von $x$ ist. Dann ist
\[
  f_n(x) = \sum_{k=0}^{n-1} \phi\bigl(T^{(k)}x\bigr)
\]
\emph{exakt additiv}: $f_{m+n}(x) = f_m(x) + f_n(T^{(m)} x)$.

\begin{tcolorbox}[ergebnis, title={Birkhoff ist \emph{direkt} anwendbar}]
Falls das Blocksystem $(X_{\mathrm{block}}, \mu_\pi, T^{(1)})$ ergodisch ist,
konvergiert
\[
  \frac{1}{n} f_n(x) = \frac{1}{n}\sum_{k=0}^{n-1}\log_2 F(B_k(x))
  \xrightarrow{n\to\infty} \int \log_2 F(B)\,d\mu_\pi = \Epi[\log_2 F(B)]
\]
für $\mu_\pi$-fast alle $x$. Kein Umweg über Subadditivität nötig ---
die Folge ist additiv.
\end{tcolorbox}

\textbf{Voraussetzung:} Ergodizität von $T^{(1)}$ auf dem Blocksystem.
Das ist die kritische offene Frage (§\,\ref{sec:voraussetzungen}).

\subsection{Kingmans Unteradditivitätssatz}

\textbf{Anwendung auf Collatz:}
Kingman wäre nötig, wenn $f_{m+n}(x) < f_m(x) + f_n(T^{(m)}x)$ echt
zutrifft. Das ist für die reine Block-Faktor-Folge \emph{nicht} der Fall
(die Blöcke hängen durch den Übergangsmechanismus zusammen, aber die
Summenstruktur bleibt additiv).

\begin{tcolorbox}[kingman, title={Wann Kingman gegenüber Birkhoff Mehrwert hat}]
Kingman ist stärker als Birkhoff genau dann, wenn die Folge $(f_n)$ echt
subadditiv ist, d.h.\ wenn die ``Kopplung'' zwischen Blöcken zur
Ungleichheit führt. Das tritt auf bei:
\begin{itemize}[nosep]
  \item \textbf{LTE-Korrekturen:} Die tatsächliche 2-adische Tiefe
    $\vii(3^{k+1}m-1)$ eines langen C-Blocks hängt vom Vorgänger $m$
    ab. Das führt zu Gedächtniseffekten.
  \item \textbf{Zufälliger Hintergrund:} Wenn man robuste obere Schranken
    für alle Orbits (statt fast alle) sucht.
  \item \textbf{Subadditive Hülle} (§\,\ref{sec:kingman-lte}): Die
    superadditive Folge $g_n(x) = \inf_{k\geq n} \frac{1}{k}f_k(x)\cdot k$
    erlaubt Aussagen über den $\limsup$.
\end{itemize}
\end{tcolorbox}

\subsection{Ruelle-Perron-Frobenius (Markov-Ketten)}

Für das endliche Markov-System auf dem mod-12-Zustandsraum ist der
präziseste Satz der Ruelle-Perron-Frobenius (RPF)-Satz:

\begin{theorem}[RPF für endliche Markov-Ketten]
\label{thm:rpf}
Sei $(Z_n)_{n\geq 0}$ eine irreduzible, aperiodische Markov-Kette auf
endlichem Zustandsraum $S$ mit stationärer Verteilung $\pi$. Sei
$\phi: S \to \R$ messbar und beschränkt. Dann gilt für $\mu_\pi$-fast alle
Startzustände $z_0$:
\[
  \frac{1}{n}\sum_{k=0}^{n-1}\phi(Z_k) \xrightarrow{n\to\infty} \Epi[\phi]
  = \sum_{s \in S} \pi(s)\,\phi(s).
\]
Außerdem gilt ein zentraler Grenzwertsatz mit Varianz
$\sigma^2 = \sum_{s}\pi(s)(\phi(s) - \Epi[\phi])^2 + 2\sum_{k\geq 1}\mathrm{Cov}_\pi(\phi(Z_0), \phi(Z_k))$.
\end{theorem}

\textbf{Warum RPF für Collatz am direktesten ist:} Der Transfer-Operator
$\mathcal{L}$ der mod-12-Markov-Kette hat eine \emph{Spektrallücke}
(der zweitgrößte Eigenwert $\lambda_2 < 1$). Das impliziert:
\begin{itemize}
  \item Exponentiell schnelle Mischung: $\|\mu_t - \pi\|_{TV} \leq C\lambda_2^t$.
  \item Konvergenz von $\frac{1}{n}\sum_{k}\phi(B_k)$ für \emph{alle}
    Startzustände, nicht nur $\mu_\pi$-fast alle.
  \item Explizite Schranke für die Mischungszeit.
\end{itemize}

\begin{center}
\renewcommand{\arraystretch}{1.4}
\begin{tabular}{>{\raggedright}p{3.2cm}
                >{\centering}p{2.6cm}
                >{\centering}p{2.6cm}
                >{\raggedright\arraybackslash}p{4.2cm}}
\toprule
\textbf{Satz} & \textbf{Folgentyp} & \textbf{Beweisstärke} & \textbf{Für Collatz} \\
\midrule
Birkhoff (1931) & additiv & $\mu$-f.ü. & Direkt, falls Ergodizität \\
Kingman (1968) & subadditiv & $\mu$-f.ü. & Für LTE-Hülle nützlich \\
RPF / Markov & additiv, endl. Zustandsr. & alle Starts & \textbf{Am direktesten} \\
\bottomrule
\end{tabular}
\end{center}

% ====================================================
\section{Prüfung aller Voraussetzungen}
\label{sec:voraussetzungen}
% ====================================================

Wir prüfen systematisch, welche Voraussetzungen im Collatz-System
erfüllt, wahrscheinlich erfüllt oder offen sind.

\subsection{Tabelle der Voraussetzungen}

\begin{center}
\renewcommand{\arraystretch}{1.5}
\begin{tabular}{>{\raggedright}p{5.0cm}
                >{\centering}p{2.2cm}
                >{\raggedright\arraybackslash}p{5.0cm}}
\toprule
\textbf{Voraussetzung} & \textbf{Status} & \textbf{Bemerkung} \\
\midrule
Maßraum $(X, \mu)$ existiert
  & \textcolor{bewiesengrün}{\bfseries ✓ Klar}
  & Haar-Maß auf $\Ztwo^*$; kompakter p-adischer Raum \\[2pt]
$T$ ist messbar
  & \textcolor{bewiesengrün}{\bfseries ✓ Klar}
  & $T$ ist stetig auf $\Ztwo^*$ (Borel-messbar) \\[2pt]
Block-Abbildung $T^{(1)}$ messbar
  & \textcolor{bewiesengrün}{\bfseries ✓ Klar}
  & Komposition stetiger Funktionen \\[2pt]
$f_n \in L^1(\mu)$
  & \textcolor{bewiesengrün}{\bfseries ✓ Wahrsch.}
  & $|\log_2 F(B)| \leq C(k+1)$; geometrische Verteilung von $k$ \\[2pt]
Endlicher Zustandsraum (mod\,12)
  & \textcolor{bewiesengrün}{\bfseries ✓ Klar}
  & $\{1,5,7,11\}$ mod\,12 \\[2pt]
Transfer-Operator irreduzibel
  & \textcolor{bewiesengrün}{\bfseries ✓ Num.}
  & Empirisch: alle 4 Zustände kommunizieren \\[2pt]
Transfer-Operator aperiodisch
  & \textcolor{bewiesengrün}{\bfseries ✓ Num.}
  & Keine periodischen Unterzyklen beobachtet \\[2pt]
Spektrallücke des Transfer-Operators
  & \textcolor{warnorange}{\bfseries $\approx$ Num.}
  & Numerisch: $\lambda_2 \approx 0{,}35$; formaler Beweis offen \\[2pt]
Ergodizität von $T$ auf $\Ztwo^*$
  & \textcolor{lückenrot}{\bfseries ✗ Offen}
  & Äquivalent zur Collatz-Vermutung (2-adisch) \\[2pt]
Maßerhaltung von $T$
  & \textcolor{lückenrot}{\bfseries ✗ Falsch}
  & $T$ ist nicht $\mu_H$-erhaltend; $T$ ist dissipativ \\
\bottomrule
\end{tabular}
\end{center}

\subsection{Die Ergodizitätsfrage}

\begin{tcolorbox}[lücke, title={Hauptproblem: Ergodizität}]
Die Ergodizität von $T$ auf $(\Ztwo^*, \mu_H)$ ist die zentrale offene
Frage. Formal: $T$ ist ergodisch $\Leftrightarrow$ jede messbare
$T$-invariante Menge hat Maß $0$ oder $1$.

\medskip
Diese Eigenschaft ist zur Collatz-Vermutung \emph{im Wesentlichen
äquivalent} über den folgenden Zusammenhang: Wenn die Collatz-Vermutung
gilt, dann hat $T$ auf dem 2-adischen Raum nur einen einzigen minimalen
Attraktor ($\{1\}$ über $\N$), was Irreduzibilität und Ergodizität impliziert.
Umgekehrt würde ein ergodisches, $\mu_H$-erhaltendes $T$ mit negativem
Lyapunov-Exponenten die Collatz-Vermutung für $\mu_H$-fast alle Startpunkte
implizieren.

\medskip
\textbf{Ausweg:} Man kann Birkhoff/RPF auf die mod-12-Markov-Kette
anwenden (endlicher Zustandsraum, klar irreduzibel und aperiodisch), ohne
die volle Ergodizität von $T$ zu benötigen.
\end{tcolorbox}

\subsection{Das Maßerhaltungsproblem und Birkhoffs Anwendbarkeit}

Da $T$ kein $\mu_H$-erhaltendes System ist, ist der klassische Birkhoff-Satz
nicht direkt auf $(\Ztwo^*, \mu_H, T)$ anwendbar. Es gibt zwei Auswege:

\begin{enumerate}[label=(\arabic*)]
  \item \textbf{Invariantes Maß suchen:} Es ist offen, ob ein $T$-invariantes
    Wahrscheinlichkeitsmaß auf $\Ztwo^*$ existiert. Tao (2019) konstruiert
    kein solches Maß explizit, arbeitet aber mit Log-Dichte-Argumenten.
  \item \textbf{Markov-Niveau:} Auf dem endlichen mod-12-Zustandsraum hat die
    induzierte Markov-Kette eine eindeutige stationäre Verteilung $\pi$.
    Das ist das ``natürliche'' invariante Maß für den Blockmechanismus.
\end{enumerate}

% ====================================================
\section{Verbindung zu Tao (2019)}
\label{sec:tao}
% ====================================================

\begin{tcolorbox}[tao, title={Taos Hauptergebnis (2019)}]
\textbf{Tao (2019):} Für jede Funktion $f: \N \to \R$ mit $f(n) \to \infty$
gilt:
\[
  \text{Dens}\bigl(\{n \in \N : \min_{k\geq 0} \mathrm{Col}^k(n) \leq f(n)\}\bigr) = 1,
\]
wobei $\text{Dens}$ die logarithmische Dichte bezeichnet.

Mit anderen Worten: Fast alle natürlichen Zahlen (bzgl.\ logarithmischer
Dichte) haben Collatz-Bahnen, die beliebig klein werden.
\end{tcolorbox}

\subsection{Impliziter Lyapunov-Exponent bei Tao}

Taos Beweis benutzt:
\begin{itemize}
  \item Fourier-Analyse auf $\Ztwo$ (2-adische Charaktere).
  \item Exponential-Summen-Abschätzungen für die Collatz-Statistik.
  \item Implizit: Die 2-adische Verteilung von $\vii(3n+1)$ wird als
    geometrisch verteilt behandelt (Heuristik der geometrisch verteilten
    Bit-Längen).
\end{itemize}

Die Kernaussage lässt sich in unserer Sprache so reformulieren:
\[
  \frac{1}{n}\sum_{k=0}^{n-1}\log_2 F(B_k(x)) < -\varepsilon
\]
gilt für Log-Dichte-$1$ viele Startzahlen $x$ und alle hinreichend großen $n$.

\subsection{Unser $\Lambda \approx -0{,}830$ als explizite Quantifizierung}

\begin{tcolorbox}[ergebnis, title={Quantifizierung von Taos Ergebnis}]
Der Lyapunov-Exponent
\[
  \boxed{\Lambda = \Epi[\log_2 F(B)]
  = \pi_E \cdot \log_2 F_E
  + \pi_A \cdot \log_2 F_A
  + \pi_B \cdot \log_2 F_B
  + \pi_C \cdot \log_2 F_C
  \approx -0{,}830}
\]
ist die explizite Quantifizierung des in Taos Beweis implizit enthaltenen
negativen Exponenten.

\medskip
\textbf{Präzisierung:} Taos Ergebnis gibt $\varepsilon > 0$ ohne expliziten
Wert. Unser $\Lambda \approx -0{,}830$ entspricht dem Erwartungswert
$\mathbb{E}_\pi[\log_2 F]$ unter der stationären Markov-Verteilung $\pi$.
Das gibt dem ``fast alle''-Ergebnis ein konkretes quantitatives Gesicht:
Jeder typische Block kontrahiert im Mittel um den Faktor $2^{-0{,}830} \approx 0{,}562$.

\medskip
\textbf{Wichtige Einschränkung:} Tao beweist Log-Dichte $1$, nicht
Haar-Maß $1$. Die Aussage gilt also für ``fast alle $n$'' im
logarithmischen Sinne, aber nicht notwendigerweise für alle $n$
(was der Collatz-Vermutung entspräche).
\end{tcolorbox}

\subsection{Warum $\Lambda \approx -0{,}830$ und nicht $\Lambda_A \approx -1{,}83$?}

Der Unterschied zwischen den zwei berechneten Werten:
\begin{itemize}
  \item $\Lambda_A \approx -1{,}83$ (aus \textit{collatz\_lte\_dichte.tex}):
    Gewichteter Mittelwert über \emph{alle} Blocklängen mit der 2-adischen
    Dichtegewichtung $P(l=k) \leq 2^{-(k+1)}$.
  \item $\Lambda \approx -0{,}830$: Mittlerer Blockfaktor unter der
    stationären Markov-Verteilung $\pi$ auf $\{E,A,B,C\}$, ohne
    Längengewichtung.
\end{itemize}
Beide Werte sind negativ --- das ist das Wesentliche. Die genauen Zahlenwerte
hängen vom Modellierungsrahmen ab.

% ====================================================
\section{Kingman für die LTE-Korrekturen: Subadditive Hülle}
\label{sec:kingman-lte}
% ====================================================

\subsection{Wo Kingman wirklich stärker ist}

Obwohl die Block-Folge additiv ist, gibt es einen Aspekt, bei dem Kingman
einen echten Mehrwert liefert: die \emph{LTE-Korrekturen}.

Für lange C-Ketten (Länge $k \geq 3$) ist der Wachstumsfaktor
$F_k = (3/2)^{k+1} \cdot 2^{-\vii(3^{k+1}m-1)}$ \emph{nicht} unabhängig
von der Vorgeschichte: Die 2-adische Tiefe $\vii(3^{k+1}m-1)$ hängt von
$m = (n_0+1)/2^{k+1}$ ab, das seinerseits durch vorherige Blöcke bestimmt wird.

\begin{tcolorbox}[kingman, title={Subadditive Hülle der LTE-korrigierten Folge}]
Definiere die \emph{LTE-korrigierte} Folge
\[
  \tilde{f}_n(x) = \text{(tatsächlicher Logarithmus des Wachstumsfaktors
  nach $n$ Blöcken mit LTE-Korrekturen)}.
\]
Falls die LTE-Korrekturen ``Gedächtnis'' einführen (was sie tun:
$\vii(3^{k+1}m-1)$ hängt von $m$ ab), dann gilt \emph{im Allgemeinen}:
\[
  \tilde{f}_{m+n}(x) \leq \tilde{f}_m(x) + \tilde{f}_n(T^{(m)}x).
\]
In Worten: Die LTE-korrigierte Folge ist subadditiv, weil lange C-Ketten
in einem zusammengesetzten Block ``besser kontrahieren'' als die Summe
zweier Einzelblöcke erwarten ließe (die Kette ``weiß'', dass $m$ groß ist).

\medskip
Kingmans Satz liefert dann:
\[
  \limsup_{n\to\infty}\frac{1}{n}\tilde{f}_n(x)
  \leq \Lambda_{\mathrm{Kingman}}
  = \inf_{n\geq 1}\frac{1}{n}\mathbb{E}[\tilde{f}_n].
\]
Falls $\Lambda_{\mathrm{Kingman}} < 0$, dann $\tilde{f}_n \to -\infty$ fast sicher.
\end{tcolorbox}

\begin{bemerkung}[Subadditivität ist nicht erzwungen]
Die LTE-Korrekturen können auch zu \emph{Superadditivität} führen, falls
die Vorgeschichte ungünstig ist (schlechter $m$-Wert). Daher ist die
Frage, ob $\tilde{f}_n$ echt subadditiv oder nur additiv ist, nicht a priori
klar. Man müsste die Korrelationsstruktur der $m$-Werte analysieren.
\end{bemerkung}

\subsection{Die superadditive Hülle als Schranke}

Eine sicherere Anwendung von Kingman ist die \emph{superadditive Hülle}:

\begin{definition}[Superadditive Hülle]
Für eine beliebige Folge $(f_n)$ definiere die superadditive Hülle:
\[
  g_n(x) = \inf_{k \geq n}\frac{1}{k}f_k(x)\cdot n.
\]
Dann ist $(g_n)$ superadditiv, und $\frac{1}{n}g_n \to f^*$ von unten.
\end{definition}

Kingmans Satz (superadditive Version) gibt:
\[
  f^*(x) = \sup_{n \geq 1}\frac{1}{n}\mathbb{E}[g_n]
  = \lim_{n\to\infty}\frac{1}{n}\mathbb{E}[g_n].
\]
Falls $f^* < 0$ (was für $\Lambda \approx -0{,}830$ plausibel ist), dann
folgt $f_n \to -\infty$ fast sicher --- unabhängig davon, ob die Folge
additiv oder subadditiv ist.

% ====================================================
\section{Was bleibt: Ergodizität als letzter offener Schritt}
\label{sec:offen}
% ====================================================

\subsection{Abhängigkeitsgraph der Schlüsselaussagen}

Die logische Struktur des Beweisversuchs:

\begin{center}
\renewcommand{\arraystretch}{1.4}
\begin{tabular}{>{\raggedright}p{5.5cm}
                >{\centering}p{2.0cm}
                >{\raggedright\arraybackslash}p{4.5cm}}
\toprule
\textbf{Aussage} & \textbf{Status} & \textbf{Benötigt} \\
\midrule
$\Lambda < 0$ (negativer Exp.) & Numerisch & Explizite Berechnung \\
Transfer-Op.\ hat Spektrallücke & Formal offen & Irreduzibilität + Aperiodizität \\
Markov-Kette (mod\,12) ergodisch & Numerisch ✓ & Irreduzibilität (beweisbar?) \\
$\frac{1}{n}f_n \to \Lambda$ fast sicher & \textcolor{warnorange}{Bedingt} & Ergodizität von $T^{(1)}$ \\
Ergodizität von $T$ auf $\Ztwo^*$ & \textcolor{lückenrot}{Offen} & Äquiv.\ zu Collatz? \\
Collatz-Vermutung (alle $n$) & \textcolor{lückenrot}{Offen} & Uniforme Schranke \\
\bottomrule
\end{tabular}
\end{center}

\subsection{Der entscheidende Lückenschluss}

\begin{tcolorbox}[lücke, title={Die letzte offene Lücke}]
Selbst wenn man akzeptiert:
\begin{enumerate}[label=(\arabic*), nosep]
  \item $\Lambda \approx -0{,}830 < 0$ (numerisch solide),
  \item Spektrallücke des Transfer-Operators (numerisch, formal offen),
  \item Quasi-Invarianz des Haar-Maßes unter $T$ (Tao),
\end{enumerate}
verbleibt der Schritt von
\[
  \text{``}\tfrac{1}{n}f_n \to \Lambda < 0 \text{ für } \mu\text{-fast alle } x\text{''}
\]
zur Collatz-Vermutung (``für \emph{alle} ungeraden $n \in \N$''):

\medskip
\textbf{Das Maß-$0$-Problem:} Birkhoff und Kingman liefern Konvergenz
``$\mu$-fast überall'', aber ein einziger Orbit könnte auf einer
$\mu$-Nullmenge liegen und trotzdem divergieren. Die Collatz-Vermutung
fordert Konvergenz auf ganz $\N$ --- einer abzählbaren Menge, die
im 2-adischen Haar-Maß ebenfalls eine Nullmenge ist.

\medskip
\textbf{Ausweg:} Man müsste zeigen, dass alle ganzzahligen Startpunkte
``typisch'' sind im Sinne der ergodischen Theorie --- d.h.\ in keiner
$T$-invarianten Nullmenge liegen. Das ist ein Uniformitätsproblem, das
durch keinen der drei Ergodensätze allein gelöst wird.
\end{tcolorbox}

\subsection{Was die ergodische Methode dennoch leistet}

\begin{tcolorbox}[ergebnis, title={Positive Leistungsbilanz der ergodischen Methode}]
\begin{enumerate}[label=(\arabic*)]
  \item \textbf{Quantitative Schranke:} $\Lambda \approx -0{,}830$ ist die
    erste explizite numerische Schranke für den mittleren Konvergenzgrad.
  \item \textbf{Taos Ergebnis präzisiert:} Log-Dichte-$1$-Aussage erhält
    konkreten Exponenten.
  \item \textbf{LTE-Barriere relativiert:} Die LTE-Worst-Case-Blöcke tragen
    exponentiell wenig zum Gesamtdrift bei (aus \textit{collatz\_lte\_dichte.tex}).
  \item \textbf{Strukturelle Kompensation bewiesen:} Nach jedem LTE-extremalen
    Block folgt erzwungene Kontraktion (Proposition aus \textit{collatz\_lte\_dichte.tex}).
  \item \textbf{RPF anwendbar auf mod-12-Ebene:} Auf dem endlichen
    Zustandsraum ist Konvergenz für alle (nicht nur fast alle) Startklassen
    garantiert, mit expliziter Mischungsrate.
\end{enumerate}
\end{tcolorbox}

% ====================================================
\section{Fazit: Bewertungsmatrix}
\label{sec:fazit}
% ====================================================

\subsection{Welcher Satz passt für Collatz?}

\begin{center}
\renewcommand{\arraystretch}{1.6}
\begin{tabular}{>{\raggedright}p{3.5cm}
                >{\centering}p{1.8cm}
                >{\centering}p{1.8cm}
                >{\centering}p{1.8cm}
                >{\raggedright\arraybackslash}p{3.2cm}}
\toprule
\textbf{Kriterium} & \textbf{Birkhoff} & \textbf{Kingman} & \textbf{RPF/Markov} & \textbf{Anmerkung} \\
\midrule
Folgentyp passend
  & additiv ✓
  & subadditiv (✓)
  & additiv ✓
  & Block-Folge additiv \\
Benötigt Maßerhaltung
  & Ja ✗
  & Ja ✗
  & Nein ✓
  & $T$ dissipativ \\
Endlicher Zustandsraum
  & Nein
  & Nein
  & Ja ✓
  & mod-12 \\
Konvergenz für alle Starts
  & Nein
  & Nein
  & Ja ✓
  & Bei Irreduzibilität \\
Ergodizität nötig
  & Ja
  & Ja
  & Nein ✓
  & RPF braucht nur Irreduz. \\
LTE-Korrekturen
  & Nicht direkt
  & Ja ✓
  & Approximativ
  & Subadditivität nützlich \\
Verbindung zu Tao
  & Mittel
  & Mittel
  & Stark ✓
  & Expliziter Exponent \\
\midrule
\textbf{Gesamteignung}
  & \textcolor{warnorange}{\bfseries Mittel}
  & \textcolor{warnorange}{\bfseries Ergänzend}
  & \textcolor{bewiesengrün}{\bfseries Am besten}
  & \\
\bottomrule
\end{tabular}
\end{center}

\subsection{Abschließende Einschätzung}

\begin{tcolorbox}[kernidee, title={Zusammenfassung: Drei Werkzeuge, ein Ziel}]
\textbf{Birkhoff} ist der ``natürliche'' Satz für die additive Block-Folge,
aber er scheitert am Maßerhaltungsproblem: $T$ ist kein erhaltendes System.
Birkhoff wäre vollständig anwendbar, falls ein $T$-invariantes Maß auf
$\Ztwo^*$ gefunden würde --- was eine offene Frage bleibt.

\medskip
\textbf{Kingman} liefert keinen direkten Vorteil für die additive Block-Folge.
Er ist jedoch das richtige Werkzeug für die LTE-korrigierte, subadditive
Hülle der Folge und für robuste obere Schranken auf dem $\limsup$.

\medskip
\textbf{Ruelle-Perron-Frobenius} ist für das Collatz-Block-System am
direktesten anwendbar: Der endliche mod-12-Zustandsraum, die numerisch
belegte Irreduzibilität und Aperiodizität, und die Spektrallücke des
Transfer-Operators ergeben zusammen den stärksten verfügbaren Konvergenzsatz
--- ohne dass volle Ergodizität von $T$ benötigt wird.

\medskip
\textbf{Offene Kernfrage:} Der Schritt von ``$\mu$-fast alle Startpunkte
konvergieren'' zu ``alle ganzzahligen Startpunkte konvergieren'' ist der
eigentliche Collatz-Inhalt. Er wird durch keinen der drei Sätze automatisch
geliefert und erfordert ein zusätzliches Uniformitätsargument.
\end{tcolorbox}

\begin{tcolorbox}[warnung, title={Ehrliche Gesamtbilanz}]
\begin{itemize}[nosep]
  \item Kingman, Birkhoff und RPF sind \textbf{leistungsstarke Werkzeuge},
    die für Collatz \textbf{bedingt anwendbar} sind.
  \item Der negative Lyapunov-Exponent $\Lambda \approx -0{,}830$
    ist robust und quantifiziert Taos Ergebnis explizit.
  \item Die Collatz-Vermutung folgt aus diesen Überlegungen \textbf{nicht}
    direkt --- die Ergodizitätslücke und das Nullmengen-Problem bleiben bestehen.
  \item Der nächste substanzielle Schritt wäre: Spektrallücke des
    Transfer-Operators formell beweisen (mod-12-Irreduzibilität ist
    endliche Kombinatorik).
\end{itemize}
\end{tcolorbox}

\subsection{Referenzen}

\begin{itemize}
  \item J.\ F.\ C.\ Kingman: \textit{The ergodic theory of subadditive stochastic
    processes}. J.\ Roy.\ Stat.\ Soc.\ Ser.\ B \textbf{30} (1968), 499--510.
  \item G.\ D.\ Birkhoff: \textit{Proof of the ergodic theorem}. Proc.\ Nat.\
    Acad.\ Sci.\ USA \textbf{17} (1931), 656--660.
  \item T.\ Tao: \textit{Almost all orbits of the Collatz map attain almost
    bounded values}. Forum Math.\ Pi \textbf{10} (2022), e12.
  \item D.\ Ruelle: \textit{Thermodynamic Formalism}. Addison-Wesley, 1978.
  \item Ergänzende Dokumente: \textit{collatz\_lte\_dichte.tex},
    \textit{collatz\_schlussartikel.tex}.
\end{itemize}

\end{document}
